\section{Discussion}\label{sec:6}

\subsection{SAVO and the rank-matched transformer}\label{sec:6.1}

The 60m exploratory ablation isolates the structural ingredient that closes the SAO$\to$QKVO gap: the $Q$-value content projection $W_c$, not multi-head routing. SAO (no $W_c$, $h$=1) sits at 272 ppl; SAVO ($W_c$, $h$=1) at 263; MH-SAO ($h$=6, no $W_c$) at 275; MH-QVC ($h$=6, $W_c$) at 266. The Transformer's $W_V$ --- a learned projection of raw $x$ --- and SAVO's $W_c$ --- a learned projection of the state$\odot$action product --- both reintroduce a content transform after content gathering, but operate on different domains: raw input vs.\ self-$Q$-value. The empirical consequence is that SAVO closes most of the SAO$\to$QKVO gap on this controlled ablation while preserving the $W_V$-free property of \QC{} at the raw-input level (no learned linear projection of raw $x$ in the content path).

\paragraph{Concurrent work and structural family.}
\QC{}~\cite{ref1} was published in March 2026, with SAVO as a four-projection variant introduced in this empirical follow-up. DeepSeek-V4~\cite{ref18} (preview release) reports an architecture in the same low-rank-routing family at trillion-parameter scale, developed concurrently. We did not draw on V4's design and the V4 preview does not cite \QC{}; the two architectures appear to have arrived at related primitives independently. The structural overlap is real and worth documenting precisely:

\begin{itemize}
  \item \textbf{Rank-$\boldsymbol{r}$ routing factorisation.} V4's Lightning Indexer queries are produced by a rank-$d_c$ down-then-up factorisation: $c^Q_t = h_t \cdot W^{DQ}$ followed by $q^I_t = c^Q_t \cdot W^{IUQ}$ (V4 eqs.\ 13--14). \QC{} uses the same down-projection structure: state $= x \cdot W_s$, action $= x \cdot W_a$, with the routing softmax operating in the rank-$r$ subspace.
  \item \textbf{Hadamard-product content compression.} V4's CSA aggregates compressed KV entries via $C^{\mathrm{Comp}}_i = \sum_j S_j \odot C_j$ (V4 eq.\ 12), an elementwise-product compression of the KV stream. SAVO's content gather is $\mathrm{content}_i = (\mathrm{state}_i \odot \mathrm{action}_i) \cdot W_c$, an elementwise product of the routing's own state and action. The elementwise product is not the same in either operand or aggregation pattern, but it occupies the same algebraic role of "low-rank content via Hadamard product."
  \item \textbf{Differences.} V4 wraps these primitives with a separate compression-weight stream $Z^a, Z^b$ producing softmax-normalised aggregation weights, top-$k$ sparse selection over the compressed entries (DSA), heavily-compressed dense attention (HCA), sliding-window KV branches, attention-sink logits, and partial RoPE on the last 64 dimensions. SAVO has none of these: it is the minimal rank-$r$ routing primitive without sparsification, hierarchical compression, or sink terms.
  \item \textbf{Empirical relationship.} V4 demonstrates that the broader low-rank-routing-with-Hadamard-content family scales to 1.6T parameters, 1M context, with 10\% of MHA's KV cache and 27\% of MHA's per-token inference FLOPs at 1M context (V4 \S1, Figure~1). The present paper provides the controlled-ablation, throughput, rank-ablation, and KV-cache analysis at 60m / 10,000-step budget for the minimal primitive in this family. The two characterisations are complementary: V4 establishes that the family is viable at scale; the present paper characterises the minimal primitive's behaviour at small scale where ablations are feasible.
\end{itemize}

We do not claim that V4 validates SAVO specifically --- the architectures differ in detail, and at 60m with 10,000 steps SAVO is $+12.33 \pm 0.87$ ppl above the rank-matched controlled ablation (4-seed paired difference, $p < 10^{-3}$, Table~\ref{tab:5}). The point of this paragraph is the narrower one: independently of \QC{}, V4 has converged on the same family of low-rank routing with Hadamard-product content gather, suggesting this family is a real point in the architecture space rather than a one-off curiosity.

\subsection{Cross-modal non-interference}\label{sec:6.2}

Unification is a structural claim; it does not automatically imply the shared parameters can absorb multiple objectives without harm. The 60m matched-text-compute observation (\S\ref{sec:5.5}) is consistent with two interpretive readings: no cross-modal transfer (auxiliary modalities do not improve text), and no cross-modal interference (auxiliary modalities do not harm text). The second reading is the non-trivial finding --- in shared-parameter multimodal architectures, interference is the default. The 180m multimodal text regression occurs in lock-step with the text-only 180m regression, so the joint training mix is not the source of the regression and the absence-of-interference observation extends through 180m.

\subsection{The world-model branch}\label{sec:6.3}

The world-model branch trains stably as a fourth concurrent objective: world MSE drops from 1.125 at initialisation to 0.287 at step 10,000 on the 60m training trajectory (Fig.~\ref{fig:4}); on held-out evaluation, the trained 180m StateEncoder reaches 0.071 (Table~\ref{tab:3}), an ${\sim}\,16\times$ reduction from initialisation in the encoded-state metric. To put 0.071 in context, we computed the predict-mean baseline post-hoc on the same trained 180m StateEncoder over the 1,000-triple held-out tail of \texttt{world\_episodes.jsonl}: $\mathrm{MSE}_{\text{predict-mean}} = 0.033$. The two numbers sit in the same band.

The reading is environmental rather than architectural. \texttt{MiniGrid-Empty-8x8-v0} is a deterministic 8$\times$8 grid in which $\mathrm{frame}_{t+1}$ differs from $\mathrm{frame}_t$ by at most one agent-cell move or rotation; the encoded next-state distribution is consequently low-variance (per-element variance ${\sim}\,0.033$). On a low-variance target distribution, a constant-mean predictor is competitive with any noisy predictor, including a trained one --- this is a known property of world-model evaluation on simple environments and is why standard world-model benchmarks~\cite{ref14,ref15} use richer environments (DMLab, Habitat, Atari) where the next-state distribution carries enough variance to separate learned dynamics from constant-mean baselines.

Within the scope of this paper --- demonstrating that the routing primitive supports a world-model objective concurrently with text/vision/audio inside a single shared backbone --- the world branch behaves as required: it trains, it converges, and it does not interfere with the other three modalities. Demonstrating world-model performance competitive with dedicated architectures (target-network EMA, predictor heads, contrastive auxiliaries, raw-pixel targets) on demanding environments is out of scope and is flagged as future work in \S\ref{sec:6.5}.

\subsection{Out-of-distribution generalisation}\label{sec:6.4}

The original OOD reading of the $W_V$-removal claim was that routing without a learned content projection of raw $x$ should generalise better off-distribution. The 60m data is consistent with this (plain SAO has the best OOD average and is best per-subset on both arxiv and pubmed). At 120m and 180m the effect does not replicate cleanly: SAVO and QKVO sit within ${\sim}$3--5 ppl of each other on arxiv, and QKVO is best on pubmed and on the OOD average at both scales. We read the 60m advantage as plausibly an under-training artifact (the $W_V$-free architecture has less capacity to overfit when the training loss has not fully converged). By 120m all variants exploit distributional content structure, and the raw-$x$ $W_V$ is at least competitive on average. Multi-head consistently hurts OOD, mirroring its in-distribution nothing-knob behaviour.

\subsection*{Summary: the primary result of this paper}

The text-LM and cross-modal results combine into a single positive structural finding: a $W_V$-free routing primitive --- with a $Q$-value content projection rather than a raw-$x$ content projection --- sits $+12.33\pm0.87$ perplexity above the rank-matched controlled-ablation transformer (4-seed paired difference, $p = 7.6 \times 10^{-4}$, Table~\ref{tab:5}) and concurrently hosts vision, audio, world-model, and four cross-field cancer-ML objectives in a single training run, without per-text-token interference. The four modalities reach their per-modality optimal scales separately (vision at 180m, audio at 60m, text at 120m) rather than uniformly, so a single scale does not dominate every modality at the 10,000-step budget. World-model performance on a chosen-simple environment is bounded by that environment's intrinsic state variance, not by the routing primitive; demanding-environment benchmarks against dedicated world-model architectures are out of scope here.

\subsection{Limitations and Open Questions}\label{sec:6.5}

\begin{enumerate}
  \item The vision encoder is lightweight (5M parameters) and not competitive with CLIP-scale encoders.
  \item The 120m multi-head variants (SAVO+MH and QKVO $h$=6) could not be run because the 120m scale config ($r$=80, $H$=640) is not divisible by $h$=6; multi-head data exists at 60m and 180m only.
  \item Both architectures regress $120m \to 180m$ at the 10,000-step training budget. The fair read is that the 180m configuration is compute-budget-starved at this step count, not that the scaling law fails structurally.
  \item The 60m text-LM head-to-head is replicated at 4 seeds (Table~\ref{tab:5}), giving measured per-architecture std 0.40--0.68 ppl --- much smaller than the architectural difference. Other comparisons (rank ablation in \S\ref{sec:5.8}, OOD subsets in \S\ref{sec:5.4}, 120m / 180m text-only scaling, full-rank MHA Phase~3) are single-seed or two-seed; a future revision should expand to $\geq$3 seeds per cell to characterise per-comparison variance more tightly.
  \item The world-model branch is evaluated on \texttt{MiniGrid-Empty-8x8-v0}, whose low encoded next-state variance puts a predict-mean baseline in the same band as the trained TransitionModel (\S\ref{sec:6.3}). Demanding environments (DMLab, Habitat, Atari) and dedicated world-model architectures (target-network EMA, predictor heads, contrastive auxiliaries, raw-pixel targets) are required to benchmark world-model learning in a stronger sense; both are out of scope here.
  \item OOD evaluation covers two scientific-prose subsets (arxiv, pubmed); dialogue and code OOD subsets and multi-seed replication are deferred.
  \item 1B-scale runs are deferred. \QC{} has not yet been empirically validated at $\geq$1B parameters.
\end{enumerate}

\subsection{What this paper does not establish}\label{sec:6.6}

We list the experiments that would be required to make stronger claims than those reported in this paper. Each item names the specific gap and the experiment that would close it.

\paragraph{Already established in this revision.}
Four items previously listed under this heading are now closed: (i) multi-seed replication of the 60m text-LM head-to-head (Table~\ref{tab:5}, 4 seeds) gives a paired SAVO$-$rank-matched difference of $+12.33 \pm 0.87$ ppl on val and $+12.16 \pm 0.54$ ppl on test, both with $p < 10^{-3}$; (ii) the SAVO compass-rank ablation across $r \in \{H/16, H/8, H/4\}$ shows $\leq 3$ ppl variation across 4$\times$ rank change (Table~\ref{tab:9}); (iii) per-attention-block throughput, FLOPs, and peak GPU memory (Table~\ref{tab:10}) show SAVO running ${\sim}6\times$ faster than full-rank standard 8-head MHA at ${\sim}5\times$ less peak GPU memory; (iv) full-rank standard 8-head MHA perplexity (Table~\ref{tab:5}, 2 seeds, batch $=2$, grad-accum$=32$ to fit 6\,GB VRAM) lands at $257.96 \pm 2.12$ val ppl. SAVO is $+5.79$ ppl \emph{above} (worse than) full-rank MHA on val. The rank-matched controlled ablation is $-6.54$ ppl \emph{below} (better than) full-rank MHA on val --- the larger 8$\times$-attn-block configuration (${\sim}$590k attn-block params/layer) is parameter-undertrained at 10,000 steps relative to the rank-matched 1$\times$-attn-block configuration (${\sim}$74k attn-block params/layer), which converges in the same budget. This is itself a substantive finding (\S\ref{sec:5.2}). The remaining open items follow.

\begin{enumerate}
  \item \textbf{10,000-step training budget.} The training budget here yields WikiText-103 val perplexity ${\sim}$250 at 60m. Comparable-parameter-count Transformer baselines in the language-modelling literature reach substantially lower perplexity at substantially longer training budgets ($10^5$--$10^6$ optimizer steps, multi-epoch); whether the SAVO/QKVO gap survives, narrows, or reverses in that converged regime is not tested in this paper.
  \item \textbf{No comparison to other parameter-efficient attention variants.} LoRA~\cite{ref25}, Linformer~\cite{ref26}, Performer~\cite{ref27}, and GQA~\cite{ref9} / MQA~\cite{ref24} are the natural reference points for any low-rank or projection-light attention proposal. None are compared empirically here. A side-by-side at matched parameter count and matched training recipe is required to position SAVO within the parameter-efficient attention literature.
  \item \textbf{World-model evaluation on a finite MDP.} \texttt{MiniGrid-Empty-8x8-v0} has ${\sim}$256 state-orientation pairs and 6 discrete actions. Whether the routing primitive supports world-model learning in the regimes that benchmark world-model architectures (Atari, DMLab, Habitat, $10^5$+ states, raw-pixel targets, dedicated training recipes with target-network EMA / predictor heads / contrastive auxiliaries) is not tested.
\end{enumerate}

These are not soft caveats. They are the experiments a reader should expect from any future revision before treating SAVO as validated against the broader attention-architecture literature. The current paper reports what it observed at the chosen scope; it does not claim what it could not test.
